\subsection{Actual affine centers, continuous model ideals and anti-continuity}
\label{subsec:native-model-continuity}

Throughout the center and model statements, $f:X\to\operatorname{Spec}R$
is an actual proper structure morphism and the prescribed real seminorm
$v$ satisfies $v(r)\le1$ for every $r\in R$. The genuine native kernel
support and the valuative model center retain their different definitions.

\begin{theorem}[Whole affine residue and structure maps]
\label{scheme:affine-native-residue-structure-maps}
\lean{EnvelopeProperty.Berkovich.SchemeValuedPoint.fromSpecResidueField_eq_affine,
EnvelopeProperty.Berkovich.SchemeValuedPoint.fromSpec_appTop_eq,
EnvelopeProperty.Berkovich.SchemeValuedPoint.fromSpec_structureMap}
\leanok
\uses{scheme:actual-residue-section-pullback,scheme:relative-affine-open-homeomorphism}
For an actual affine open $U$ and $x\in U$, its genuine section evaluation
$\Gamma(X,U)\to\kappa(x)$ gives the entire original morphism
$\operatorname{Spec}\kappa(x)\to X$ after the actual affine inclusion.
The coefficient action derived from the structure morphism presents the
entire affine structure map, not merely its numerical seminorm values.
\end{theorem}
\begin{proof}
The actual affine germ morphism followed by the residue map is the
section evaluation homomorphism. The genuine affine presentation of
$\operatorname{Spec}\mathcal O_{X,x}\to X$ therefore gives the residue
formula on whole scheme morphisms. Actual global-section pullback along
the affine inclusion is ordinary restriction followed by the true
$\Gamma\operatorname{Spec}$ isomorphism. Apply the contravariant
composition formula and naturality of that isomorphism to the actual
structure morphism. Its resulting ring homomorphism is exactly the
previously derived section coefficient map, proving the whole structure
formula by the true morphism-to-affine-spectrum equivalence.
\end{proof}

\begin{theorem}[Actual bounded affine criterion for proper centers]
\label{scheme:proper-center-bounded-affine-criterion}
\lean{EnvelopeProperty.Berkovich.RelativeSchemeValuedPoint.boundedAffineSectionMap,
EnvelopeProperty.Berkovich.RelativeSchemeValuedPoint.boundedAffineSectionMap_base,
EnvelopeProperty.Berkovich.RelativeSchemeValuedPoint.boundedAffineLift,
EnvelopeProperty.Berkovich.RelativeSchemeValuedPoint.boundedAffineLift_generic,
EnvelopeProperty.Berkovich.RelativeSchemeValuedPoint.boundedAffineLift_base,
EnvelopeProperty.Berkovich.RelativeSchemeValuedPoint.boundedAffineLift_eq_centerLift,
EnvelopeProperty.Berkovich.RelativeSchemeValuedPoint.center_mem_of_bounded_affine,
EnvelopeProperty.Berkovich.RelativeSchemeValuedPoint.center_mem_of_mem_boundedAffineImage,
EnvelopeProperty.Berkovich.RelativeSchemeValuedPoint.mem_boundedAffineImage_iff_center_mem,
EnvelopeProperty.Berkovich.RelativeSchemeValuedPoint.boundedAffineImage_eq_center_preimage}
\leanok
\uses{scheme:affine-native-residue-structure-maps,scheme:proper-native-center,
scheme:proper-center-local-bounds,scheme:compact-bounded-native-affine-images}
A native point supported in an actual affine open $U$ has its true proper
center in $U$ if and only if its evaluation on every actual section of
$U$ is at most one. Consequently, the genuine bounded affine chart image
is exactly $\operatorname{center}^{-1}(U)$.
\end{theorem}
\begin{proof}
Bounded section values put the entire actual section-evaluation
homomorphism into the original residue-field valuation ring.
Genuine germ restriction proves equality with the original coefficient
homomorphism inside this valuation ring; equality of norms alone is
not used. Its spectrum morphism followed by the actual affine chart
inclusion gives a genuine whole valuation-ring lift into $X$.
The preceding affine residue and structure identities prove both actual
valuative triangles. Proper separated uniqueness identifies this whole
lift with the previously constructed model lift. Its closed-point image
lies in $U$, proving the required center membership. The converse is
the actual section bound already proved from model lift pullback.

The genuine relative affine homeomorphism preserves every section value.
Thus a bounded chart point supplies exactly the required bound, and an
actual center neighborhood supplies a bounded chart representative.
These two directions identify the entire bounded affine image with the
true center inverse image. The argument assumes neither the center
criterion nor an unrelated bounded affine coefficient action.
\end{proof}

\begin{theorem}[Closed affine pieces and their continuous model functions]
\label{scheme:closed-bounded-affine-model-functions}
\lean{EnvelopeProperty.Berkovich.RelativeSchemeValuedPoint.isClosed_boundedAffineImage,
EnvelopeProperty.Berkovich.RelativeSchemeValuedPoint.isClosed_center_preimage_affine,
EnvelopeProperty.Berkovich.RelativeSchemeValuedPoint.continuousOn_modelIdealNorm_boundedAffineImage,
EnvelopeProperty.Berkovich.RelativeSchemeValuedPoint.continuousOn_modelIdealLog_boundedAffineImage}
\leanok
\uses{scheme:separated-native-hausdorff,scheme:compact-bounded-native-affine-images,
scheme:proper-center-bounded-affine-criterion,scheme:native-model-ideal-generators}
The genuine bounded affine image is closed for every separated structure
morphism to the affine base. On a proper model it is the true center
inverse image. If an actual ideal-sheaf-data ideal $I(U)$ has specified
finite generators, the intrinsic model norm and logarithm are continuous
on this closed affine image.
\end{theorem}
\begin{proof}
The independently proved native Hausdorff property makes each genuinely
compact bounded affine image closed. The actual center criterion gives
its center interpretation. Every point of this piece is an actual
center-neighborhood point, so the previously proved intrinsic model norm
formula equals the finite maximum of its generator section values.
The genuine relative affine homeomorphism is continuous on the piece
and preserves all sections. Continuous coordinate evaluation, finite
maxima and the continuous inclusion into the extended nonnegative reals
prove norm continuity on that piece. The actual continuous extended
logarithm proves logarithmic continuity, including norm zero and
negative infinity. No globally continuous center map is assumed.
\end{proof}

\begin{theorem}[Global continuous genuine model ideal functions]
\label{scheme:continuous-native-model-ideals}
\lean{EnvelopeProperty.Berkovich.RelativeSchemeValuedPoint.continuous_modelIdealNorm_of_finite_affine_cover,
EnvelopeProperty.Berkovich.RelativeSchemeValuedPoint.continuous_modelIdealNorm_of_fg,
EnvelopeProperty.Berkovich.RelativeSchemeValuedPoint.continuous_modelIdealLog_of_fg}
\leanok
\uses{scheme:closed-bounded-affine-model-functions,scheme:proper-native-compactness}
A genuine finite affine scheme cover on which the actual ideals have
finite generators gives a globally continuous intrinsic model norm.
In particular, if $I(U)$ is finitely generated for every actual affine
open, its intrinsic model norm and logarithm are globally continuous.
No noetherian hypothesis is needed when this finite generation is supplied.
\end{theorem}
\begin{proof}
The actual scheme affine cover yields a cover of the whole native space
by its genuine bounded affine images, because every true proper center
lies in one member. The pieces are compact and closed, and the preceding
result proves continuity of the same intrinsic function on each piece.
A finite family is locally finite, so the genuine closed pasting theorem
proves global continuity. The intrinsic stalk definition ensures the
functions already agree on intersections, without chosen-generator
compatibility assumptions.

Proper quasi-compactness over the actual compact affine base supplies a
real finite affine cover. Supplied finite generation chooses actual finite
generators on those section rings. Apply the finite-cover theorem and
then the continuous extended logarithm. No finite-type hypothesis on an
arbitrary open section ring, analytic structure sheaf, or continuity
certificate is introduced.
\end{proof}

\begin{theorem}[Actual center anti-continuity with explicit compact-open scope]
\label{scheme:proper-center-anticontinuity}
\lean{EnvelopeProperty.Berkovich.RelativeSchemeValuedPoint.isClosed_center_preimage_of_isCompact,
EnvelopeProperty.Berkovich.RelativeSchemeValuedPoint.isClosed_center_preimage_of_noetherian,
EnvelopeProperty.Berkovich.RelativeSchemeValuedPoint.isOpen_center_preimage_of_isClosed}
\leanok
\uses{scheme:closed-bounded-affine-model-functions}
The true center inverse image of every quasi-compact actual scheme open
is closed. If the underlying proper model is topologically noetherian,
this holds for every scheme open; equivalently, every closed subset has
an open center inverse image. No unrestricted anti-continuity assertion
for arbitrary scheme opens is made.
\end{theorem}
\begin{proof}
The genuine affine basis expresses each quasi-compact scheme open as a
finite union of actual affine opens. Their center inverse images are
closed by the preceding result. Preimages commute with the finite union,
so the required inverse image is closed. On a topologically noetherian
scheme every open is quasi-compact. Apply this statement to the open
complement of a closed subset and take complements to obtain the genuine
anti-continuity formulation. This uses topological noetherianity, not an
unproved finite generator assertion for every section ring.
\end{proof}

\begin{theorem}[Open center inverse images of genuine finitely generated ideal supports]
\label{scheme:fg-ideal-support-center-open}
\lean{EnvelopeProperty.Berkovich.RelativeSchemeValuedPoint.modelIdealNorm_lt_one_iff_of_fg,
EnvelopeProperty.Berkovich.RelativeSchemeValuedPoint.isOpen_center_preimage_idealSupport_of_fg}
\leanok
\uses{scheme:model-ideal-strict-center-criterion,scheme:continuous-native-model-ideals}
If $I(U)$ is finitely generated on every actual affine open, then
$\|I\|(y)<1$ holds exactly when the true model center of $y$ is in the
actual ideal-sheaf-data support. This support has an open center inverse
image even without a noetherian hypothesis on the model.
\end{theorem}
\begin{proof}
Choose a real affine neighborhood of the true center and actual finite
generators of its ideal. The independently proved finite-generator strict
criterion identifies the intrinsic norm inequality with genuine ideal
support membership. Global model norm continuity then makes its strict
sublevel set below one open. The finite-generation hypothesis is essential
for this converse: an infinite supremum of values strictly below one need
not have a uniform strict gap. These statements concern the true model
center, whereas vanishing of the model norm continues to detect the
original native kernel support as proved earlier.
\end{proof}
